\FloatBarrier
\label{sec:ScintData}
This chapter presents the atmospheric muon data recorded with the scintillator detector and its readout system (sec.\ \ref{sec:Scint}). The required calibration measurements of the readout system are also presented in this section.

\FloatBarrier
\section{Discriminator threshold calibration}
%\begin{itemize}
%    \item FPGA-internal comparators (on FPGA silicon die) are subject to production uncertainties
%    \item Unfolds in voltage bias on comparator switching point
%    \item Scan input voltage range while keeping other comparator input constant to detemine switching point of comparator (Metastability calibration)
%    \item Account for measured biases when setting thresolds
%\end{itemize}
\label{sec:ScintData-VoltageBias}
\FloatBarrier
\paragraph{Comparator bias}
The \ac{FPGA}-internal comparators used for discrimination of the analog \ac{SiPM} signals have individual bias voltages (sec.\ \ref{sec:Scint-ReadoutBoard-FPGASampling}). An ideal comparator would switch its output state at the point where both input voltages are equal: $V_+ = V_-$. For the non-ideal comparator, the switching point is shifted by the bias voltage: $V_+ + V_\text{bias} = V_-$. As explained, the bias voltage is expected to be as high as $\pm 35 \;\si{\milli\volt}$ \cite{UltraScale-FPGA-SelectIO-Resources}. As the discriminator threshold lies in the range of $\mathcal{O}(\SI{10}{\milli\volt})$, a calibration of the bias voltage for each individual comparator is essential.
\FloatBarrier
\paragraph{Measurement with function generator}
There are several ways to measure the real switching point of the comparator and extract the bias. The established method is to use a well-defined triangular pulse generated by a function generator for this purpose. By measuring the timestamps of the comparator switching points one can calculate the bias \cite{FPGA-MVT-Xi,FPGA-MVT-Eliseev}. The advantage of this method is that one measures the rising and falling edges of the comparator separately, i.e.\ one can conclude whether the comparator shows a hysteresis effect. However, the conducted measurements indicate no significant hysteresis.
\FloatBarrier
\paragraph{Measurement with comparator metastability}
A novel approach for the measurement of the comparator bias was introduced by our group \cite{MSc-Ehlert}, which eliminates the need for an external function generator. Instead, the method only requires an adjustable voltage source (an \ac{ADC}) at the signal input of the comparator. This significantly reduces the complexity of the calibration setup. In the case of the developed readout board (sec.\ \ref{sec:Scint-ReadoutBoard}), this functionality is already included on the \ac{PCB} itself.
\newline One input of the comparator is set to the desired discriminator threshold voltage: $V_-=V_\text{thres}$. The voltage of the other input is varied in the range $V_+ \in [V_\text{thres}-\SI{50}{\milli\volt},\; V_\text{thres}+\SI{50}{\milli\volt} ]$ with a fine step size. For each step, the count rate of the comparator switching events is measured. This is enabled by counting registers inside the \ac{FPGA} logic which can be read out and reset to zero after changing the voltage level. At the switching point of the comparator, the output is undefined. For an ideal device, one would expect a delta peak. For the real comparator, one would expect a peak in the count rate with a finite width, as electronic noise leads to small fluctuations in the comparator input voltages. By estimating the position of the peak, the bias voltage $V_\text{bias}$ can be estimated. Fig.\ \ref{fig:ScintData-VoltageBias-Scheme} illustrated the described calibration method exploiting the comparator metastability.
\begin{figure}[htb]
    \centering
    \includegraphics[width=0.8\textwidth]{Scint-VoltageBias-Scheme.pdf}
    \caption{Illustration of the described metastability calibration method to measure the comparator bias voltage \cite[adapted]{MSc-Ehlert}.}
    \label{fig:ScintData-VoltageBias-Scheme}
\end{figure}
\FloatBarrier
\paragraph{Calibration results}
The voltage scans are performed for all input channels of the \ac{FPGA} and the switching rate histogram is analyzed. Fig.\ \ref{fig:ScintData-VoltageBias-SwitchingCount} shows an example histogram of the switching rate depending on the applied \ac{DAC} voltage. The obtained bias voltages of the comparators are shown in blue. They lie in the range \SI{-10}{\milli\volt} to \SI{15}{\milli\volt}, which is compatible with the mentioned maximum range of $\pm 35 \;\si{\milli\volt}$ from the \ac{FPGA} vendor datasheet \cite{UltraScale-FPGA-SelectIO-Resources}. 
\begin{figure}[htb]
    \centering % presented histogram: ca3 zynq51 input4 step0
    \includegraphics[width=0.8\textwidth]{scint_thres_calib/single_count_hist.pdf}
    \caption{Example switching rate histogram for a single comparator of the \ac{FPGA}. The weighted mean of the distribution is marked in red. Initially it deviates from the threshold voltage (gray) applied to the other input of the comparator, because of the nonzero bias voltage.}
    \label{fig:ScintData-VoltageBias-SwitchingCount}
\end{figure}
\newline The bias voltage is calculated as difference between the switching rate peak position $\left\langle V_\text{DAC} \right\rangle$ and the expected peak position $V_\text{thres}$:
\begin{equation}
    V_\text{bias} = \left\langle V_\text{DAC} \right\rangle - V_\text{thres}.
    \label{eq:Scint-VoltageBias-BiasCalculation}
\end{equation}
 The result for the measured bias voltage for all 32 inputs of a readout board is presented in fig.\ \ref{fig:ScintData-VoltageBias-CalibResult} with the blue data points.
\begin{figure}[htb]
    \centering % presented histogram: ca3 zynq51
    \includegraphics[width=0.8\textwidth]{scint_thres_calib/bias_scatter.pdf}
    \caption{Comparison of measured comparator bias before (blue) and after (red) the calibration for all 32 inputs of one readout board. The behavior of an ideal comparator ($V_\text{bias}=\SI{0}{\milli\volt}$) is marked in gray.}
    \label{fig:ScintData-VoltageBias-CalibResult}
\end{figure}
\newline After the measurement, the comparator threshold voltages are adjusted accordingly. To account for the obtained bias voltage, the new threshold voltage is selected as
\begin{equation}
    V_\text{thres,new} = V_\text{thres,initial} - V_\text{bias}.
    \label{fig:ScintData-VoltageBias-NewThresold}
\end{equation}
The procedure is repeated multiple times, in order to iteratively improve the precision of the calibration. After setting the final threshold voltage, the bias voltage is measured one more time to validate the outcome of the calibration. The obtained new bias voltages are shown as the red data points in fig.\ \ref{fig:ScintData-VoltageBias-CalibResult}. As expected, the bias voltage could be mostly removed with this procedure, while few channels show a remaining offset of $\mathcal{O}(\SI{1}{\milli\volt})$.
% The obtained bias voltages seem to be mostly independent of the measurement time and the temperature. However, the bias depends strongly on the applied voltages to the comparator, i.e.\ the selected threshold voltage. 

\FloatBarrier
\section{TDC timing calibration}
%\begin{itemize}
%    \item Bank-specific calibration of deserializers on each power-cycle leads to non-deterministic signal delays
%    \item Use simultaneous testpulses to measure these delays
%    \item Use delay module in firmware to correct for these effects already on the device
%    \item Mention ringing of testpulse signal in circuit, but explain that no negative effect on calibration
%\end{itemize}
\label{sec:ScintData-TimingBias}
As explained in sec.\ \ref{sec:Scint-ReadoutFirmware-SamplingTDC}, the \ac{TDC} logic of the readout board is realized in the \ac{FPGA} firmware by utilizing the deserializer primitives of the chip. On startup of the \ac{FPGA}, the deserializers are internally calibrated, individually for each component. This individual calibration leads to non-deterministic time offsets between different groups of inputs. 
\FloatBarrier
\paragraph{Calibration with testpulses}
The readout board offers the possibility to transmit simultaneous testpulses to all input channels (sec.\ \ref{fig:Scint-ReadoutBoard-AnalogCircuit}). Similar to the \ac{DT} timing calibration with testpulses (sec.\ \ref{sec:DTData-Testpulse}), it is possible to extract timing offsets between the input channels by measuring the timestamps of the testpulse hits. A histogram of the recorded testpulse hit timestamps for one channel is presented in fig.\ \ref{fig:ScintData-TimingBias-TPTimestamp}.
\begin{figure}[htb]
    \centering % presented histogram: 
    \includegraphics[width=0.8\textwidth]{scint_timing_calib/single_timing_hist.pdf}
    \caption{Example histogram of testpulse arrival timing for one input channel of the readout board.}
    \label{fig:ScintData-TimingBias-TPTimestamp}
\end{figure}
\FloatBarrier
\paragraph{Programming channel delays}
\begin{figure}[htb]
    \centering % presented histogram: 
    \includegraphics[width=0.8\textwidth]{scint_timing_calib/timing_scatter.pdf}
    \caption{.}
    \label{fig:ScintData-TimingBias-CalibResult}
\end{figure}

\FloatBarrier
\section{SiPM noise scan}
%\begin{itemize}
%    \item Scan of SiPM rates for different thresolds with developed readout system
%    \item Show resulting curve
%    \item Discuss visible exponential (expected) behavior
%    \item Aliasing effects (too coarse scanning steps) hide expected steps in spectrum
%\end{itemize}
\label{sec:ScintData-Noise}

\FloatBarrier
\section{Signal ringing and crosstalk}
\label{sec:ScintData-RingingCrosstalk}

\FloatBarrier
\section{SiPM hits}
%\begin{itemize}
%    \item Show single SiPM occupancies
%    \item Discuss observed secondary pulses: Ringing in circuit or SiPM after-pulses
%\end{itemize}
\label{sec:ScintData-Hits}

\FloatBarrier
\section{SiPM coincidences}
%\begin{itemize}
%    \item Coincidence of two SiPMs of same scintillator strip
%    \item Discuss strip rate
%    \item Show time difference between both SiPMs, discuss width / time localization
%    \item Compare to expected fake coincidence rate
%\end{itemize}
\label{sec:ScintData-Strips}

\FloatBarrier
\section{Scintillator strip coincidences}
%\begin{itemize}
%    \item Coincidence of two strips of two detector layers (create "pixels")
%    \item Discuss pixel rate
%    \item Show time difference between both strips
%    \item Compare to expected fake coincidence rate
%\end{itemize}
\label{sec:ScintData-Pixels}





